\documentclass[../main.tex]{subfiles}
\begin{document}

\section{Jacobi Method}
\begin{center}
    \begin{minipage}{0.95\textwidth}
    \raggedright
    \begin{itemize*}[label=\textbullet, itemjoin=\hspace{2em}]
        \item \textbf{Linear system:} $Ax=b$
        \item \textbf{Residual:} $r=b-Ax_k$
        \item \textbf{Error:} $e=x-x_k$
    \end{itemize*} \\
    \vspace{1em}
    \begin{itemize*}[label=\textbullet, itemjoin=\hspace{2em}]
        \item \textbf{Preconditioner:} $P\approx A$
    \end{itemize*}
    \end{minipage}
\end{center}

The Jacobi method is a simple iterative method for solving linear systems of equations. It is a first-order method, and is given by the following update rule:

$$
x^{k+1} = D^{-1} (b - (L + U)x^k)
$$

where $D$ is the diagonal of the matrix $A$, $L$ is the lower triangular part of $A$, and $U$ is the upper triangular part of $A$.

Source: 
\begin{itemize}
    \item \textit{A Multigrid Tutorial} by Briggs, et al. (2000)
    \item https://yaningliucudenver.github.io/Numerical-Analysis-I/bookchapter3-2.html
\end{itemize}

\end{document}
